% ════════════════════════════════════════════════════════════════
\section{Wykład 10}
\subsection{Problem szeregowania zadań (cd.)}
% ════════════════════════════════════════════════════════════════

Zbiór $A \subseteq S$ nazywamy \underline{wolnym od kar}, jeśli istnieje harmonogram zbioru $A$, w którym wszystkie zadania są terminowe.

Dla $t \in \{0, 1, \ldots, n\}$ niech
\begin{flalign*}
	N_t(A) = |\{z_i \in A : d_i \leq t\}| && &&
\end{flalign*}
-- liczba zadań w $A$ o deadlinie $\leq t$.

\begin{lemma} \label{lemma:wolny-od-kar}
	Niech $A \subseteq S$. Następujące warunki są równoważne:
	\begin{enumerate}
		\item $A$ jest wolny od kar.
		\item $N_t(A) \leq t$ dla każdego $t \in \{1, 2, \ldots, n\}$.
		\item Jeśli zadania z $A$ są posortowane w kolejności niemalejącej ze względu na deadliny, to żadne zadanie nie jest spóźnione.
	\end{enumerate}
\end{lemma}

\begin{theorem}
	Niech $S$ -- zbiór zadań o jednostkowym czasie wykonania, a $\mathcal{A}$ -- rodzina zbiorów wolnych od kar. Wtedy $(S, \mathcal{A})$ jest matroidem.
\end{theorem}
\begin{proof}
	Sprawdzamy własności z definicji matroidu:
	\begin{enumerate}
		\item Oczywiste. $\checkmark$
		\item Dla każdego $i \in [n]$ zachodzi $\{z_i\} \in \mathcal{A}$, więc $\mathcal{A} \neq \emptyset$. $\checkmark$
		\item Niech $A \in \mathcal{A}$, $B \subseteq A$. Istnieje harmonogram zbioru $A$, w którym wszystkie zadania są terminowe. Harmonogram zbioru $B$ powstały z usunięcia z tego harmonogramu zadań z $A \setminus B$ jest harmonogramem zbioru $B$, w którym wszystkie zadania są terminowe. Zatem $B \in \mathcal{A}$.
		\item Niech $A, B \in \mathcal{A}$, $|B| > |A|$.\\
		Niech $k \in \{0, 1, \ldots, n\}$ będzie największą liczbą $t$ taką, że $N_t(B) \leq N_t(A)$.\\
		$k$ jest dobrze zdefiniowane, bo $N_0(B) = 0 = N_0(A)$.\\
		Mamy $N_n(B) = |B| > |A| = N_n(A)$, więc $k < n$.\\
		Dla $j \in \{k+1, k+2, \ldots, n\}$ mamy $N_j(B) > N_j(A)$, więc $N_j(B) \geq N_j(A) + 1$.
		\begin{flalign*}
			\left.
			\begin{array}{l}
				N_{k+1}(B) > N_{k+1}(A) \\
				N_k(B) \leq N_k(A)
			\end{array}
			\right\} \implies \text{istnieje zadanie } z \in B \setminus A \text{ o deadlinie } = k+1 && &&
		\end{flalign*}
		Niech $A' = A \cup \{z\}$. Pokażemy, że $A'$ jest wolny od kar, używając lematu \ref{lemma:wolny-od-kar}.
		\begin{itemize}[nosep]
			\item Dla $t \in \{1, 2, \ldots, k\}$:\quad $N_t(A') = N_t(A) \leq t$, bo $A$ jest wolny od kar.
			\item Dla $t \in \{k+1, k+2, \ldots, n\}$:\quad $N_t(A') = N_t(A) + 1 \leq N_t(B) \leq t$, bo $B$ jest wolny od kar.
		\end{itemize}
		Zatem $A'$ jest wolny od kar.
	\end{enumerate}
	Pokazaliśmy, że $(S, \mathcal{A})$ jest matroidem.
\end{proof}

Możemy do rozwiązania problemu szeregowania zadań użyć algorytmu Greedy (Algorytm \ref{alg:greedy}) dla matroidu $(S, \mathcal{A})$ maksymalizującego $\sum_{z_i \in A} w_i$ (czyli minimalizującego sumę kar zadań spóźnionych).

Z twierdzenia Rado, Edmonds algorytm Greedy użyty do tego problemu zwraca rozwiązanie optymalne. Warunek z linii 5 algorytmu Greedy (tzn. czy $A \cup \{x_i\} \in \mathcal{A}$) można sprawdzić w czasie $\BO(n)$ -- lemat \ref{lemma:wolny-od-kar}, warunek 3.

Złożoność czasowa algorytmu wynosi:
\begin{flalign*}
	\BO(n \log n + n \cdot n) = \BO(n^2) && &&
\end{flalign*}

\begin{example}{Szeregowanie zadań z deadlinami}
\begin{table}[H]
	\centering
	\renewcommand{\arraystretch}{1.2}
	\begin{NiceTabular}{c|cccccccc}[margin,hvlines]
		$i$   & 1  & 2  & 3  & 4  & 5  & 6  & 7  & 8  \\
		\hline
		$d_i$ & 3  & 2  & 2  & 3  & 1  & 4  & 6  & 5  \\
		$w_i$ & 80 & 70 & 60 & 50 & 40 & 30 & 20 & 10
	\end{NiceTabular}
\end{table}

Zadania są już posortowane w porządku nierosnącym względem wagi ($w_1 \geq w_2 \geq \ldots \geq w_8$). Kolejne iteracje algorytmu Greedy (poniżej $A$ wypisane w kolejności niemalejącej po deadlinach):
\begin{enumerate}[nosep]
	\item $A = (z_1)$
	\item $A = (z_2, z_1)$
	\item $A = (z_2, z_3, z_1)$
	\item $A = (z_2, z_3, z_1)$ \quad ($z_4$ nie dodajemy, bo $N_3(A \cup \{z_4\}) = 4 > 3$)
	\item $A = (z_2, z_3, z_1)$ \quad ($z_5$ nie dodajemy, bo $N_2(A \cup \{z_5\}) = 3 > 2$)
	\item $A = (z_2, z_3, z_1, z_6)$
	\item $A = (z_2, z_3, z_1, z_6, z_7)$
	\item $A = (z_2, z_3, z_1, z_6, z_8, z_7)$
\end{enumerate}

Harmonogram (zadania terminowe posortowane po deadlinach, potem zadania spóźnione):
\begin{flalign*}
	(z_2,\ z_3,\ z_1,\ z_6,\ z_8,\ z_7,\ z_4,\ z_5) && &&
\end{flalign*}
\begin{flalign*}
	\text{Suma kar} = w_4 + w_5 = 50 + 40 = 90 && &&
\end{flalign*}
\end{example}

% ════════════════════════════════════════════════════════════════
\subsection{Algorytmy aproksymacyjne}
% ════════════════════════════════════════════════════════════════

Algorytmy aproksymacyjne stosuje się do problemów, dla których jest mała szansa na znalezienie algorytmu dokładnego (np. wielomianowego). Zamiast szukać rozwiązania optymalnego, szukamy rozwiązania ,,prawie optymalnego''.

Niech $\Pi$ -- problem optymalizacyjny znalezienia maksimum lub minimum pewnej funkcji celu $f$.
\begin{itemize}[nosep]
	\item $D_\Pi$ -- zbiór instancji problemu $\Pi$.
	\item $I \in D_\Pi$ -- instancja $\Pi$.
	\item $S(I)$ -- zbiór rozwiązań dopuszczalnych dla instancji $I$.
	\item $f: \bigcup_{I \in D_\Pi} S(I) \to \R$ -- funkcja celu. Możemy założyć, że
	\begin{flalign*}
		\forall I \in D_\Pi\ \forall s \in S(I)\ f(s) > 0 && &&
	\end{flalign*}
\end{itemize}

Niech $c^*(I)$ oznacza wartość rozwiązania optymalnego instancji $I$, $c^*(I) > 0$.\\
$A$ -- algorytm aproksymacyjny rozwiązujący $\Pi$.\\
$c(I)$ -- wartość rozwiązania znalezionego przez algorytm $A$ dla instancji $I$, $c(I) > 0$.

\begin{definition}
	Mówimy, że algorytm $A$ ma \underline{współczynnik aproksymacji} $\rho(n)$, jeśli dla każdej instancji $I \in D_\Pi$ rozmiaru $n$ zachodzi
	\begin{flalign*}
		\max\left\{\frac{c(I)}{c^*(I)},\ \frac{c^*(I)}{c(I)}\right\} \leq \rho(n) && &&
	\end{flalign*}
	Wtedy $A$ nazywamy algorytmem $\rho(n)$-aproksymacyjnym.
\end{definition}

Zauważmy, że:
\begin{flalign*}
	\left.
	\begin{array}{ll}
		\text{dla problemu maksymalizacji:} & 0 < c(I) \leq c^*(I) \\[0.5ex]
		\text{dla problemu minimalizacji:}  & 0 < c^*(I) \leq c(I)
	\end{array}
	\right\} \implies \max\left\{\frac{c(I)}{c^*(I)},\ \frac{c^*(I)}{c(I)}\right\} \geq 1 && &&
\end{flalign*}

Niech $\rho = \sup_{n \in \N}\{\rho(n)\}$. Wtedy mówimy, że \underline{współczynnik aproksymacji} algorytmu $A$ wynosi $\rho$ i $A$ nazywamy algorytmem $\rho$-aproksymacyjnym.

% ════════════════════════════════════════════════════════════════
\subsection{Problem pokrycia wierzchołkowego}
% ════════════════════════════════════════════════════════════════

\begin{definition}
	Niech $G = (V, E)$ -- graf. \underline{Pokryciem wierzchołkowym} (Vertex Cover, VC) grafu $G$ nazywamy zbiór $W \subseteq V$ taki, że dla każdej krawędzi $e \in E$ zachodzi $W \cap e \neq \emptyset$ -- czyli każda krawędź $G$ ma co najmniej jeden koniec w $W$.
\end{definition}

\paragraph{Problem pokrycia wierzchołkowego}
\begin{itemize}[nosep]
	\item Instancja: graf $G$.
	\item Problem: znaleźć pokrycie wierzchołkowe grafu $G$ o minimalnej liczbie wierzchołków.
\end{itemize}

\begin{example}{Pokrycie wierzchołkowe}
\begin{figure}[H]
	\centering
	\begin{tikzpicture}[every node/.style={circle, draw, minimum size=0.6cm, inner sep=1pt}]
		\node[fill=accent!25] (a) at (0, 2) {$a$};
		\node              (b) at (2, 2) {$b$};
		\node[fill=accent!25] (e) at (1, 1) {$e$};
		\node              (c) at (0, 0) {$c$};
		\node[fill=accent!25] (d) at (2, 0) {$d$};
		\node[fill=accent!25] (f) at (4, 1) {$f$};
		\draw (a) -- (b);
		\draw (a) -- (c);
		\draw (a) -- (e);
		\draw (b) -- (d);
		\draw (b) -- (e);
		\draw (c) -- (d);
		\draw (d) -- (e);
		\draw (d) -- (f);
	\end{tikzpicture}
\end{figure}
$W = \{a, e, d, f\}$ jest pokryciem wierzchołkowym powyższego grafu $G$ (zaznaczone wyróżnione wierzchołki).
\end{example}

\paragraph{Algorytm Approx Vertex Cover}

\begin{alg}{Approx Vertex Cover}{alg:avc}
	\FunctionNN{AVC}{$G$}
	\State $C \Assign \emptyset$ \Comment{$C$ -- rozwiązanie}
	\State $F \Assign E(G)$ \Comment{$F$ -- zbiór dotychczas niepokrytych krawędzi}
	\While{$F \neq \emptyset$}
	\State wybierz dowolną krawędź $uv \in F$
	\State $C \Assign C \cup \{u, v\}$
	\State usuń z $F$ wszystkie krawędzie incydentne z $u$ lub z $v$
	\EndWhile
	\State \Return $C$
	\EndFunction
\end{alg}

\begin{example}{Algorytm AVC -- przykład}
Rozważmy graf z poprzedniego przykładu.

Pierwszy przebieg algorytmu może wybrać kolejno krawędzie $ab$ i $cd$:
\begin{itemize}[nosep]
	\item Po wybraniu $ab$ usuwamy z $F$ krawędzie $ab, ac, ae, bd, be$, zostają $cd, de, df$.
	\item Po wybraniu $cd$ usuwamy z $F$ krawędzie $cd, de, df$ -- $F = \emptyset$.
\end{itemize}
Algorytm zwraca $C = \{a, b, c, d\}$.

Inny przebieg algorytmu może wybrać kolejno krawędzie $df$, $be$ i $ac$:
\begin{itemize}[nosep]
	\item Po wybraniu $df$ usuwamy z $F$ krawędzie $df, de, bd, cd$.
	\item Po wybraniu $be$ usuwamy z $F$ krawędzie $be, ae, ab$.
	\item Po wybraniu $ac$ usuwamy ostatnią krawędź $ac$ -- $F = \emptyset$.
\end{itemize}
Algorytm zwraca $C = \{d, f, e, b, a, c\}$.
\end{example}